\section{Extended record: migrated arguments and finite evidence}
\label{long269:long:extended}

This part of the record holds material that the short note compressed or moved.
Nothing here is an additional result; every statement is either an auxiliary
form of a theorem of the note, a historical or bibliographic detail, or the
exact source inventory.  The canonical references are the note's own labels.

\subsection{The historical and catalogue record}\label{long269:long:history}

In the primary 1988 source Erd\H{o}s states the infinite-$P$ assertion as a
simple exercise and presents persistence for a finite number of primes greater
than one as a probable extension, and not as a theorem
\cite[p.~106]{erdos1988}.  The current open status is supplied by the
catalogue \cite{erdosproblems}, and is not inferred from that conjectural
wording.  The 1974 letter writes ``given primes $p_1,\ldots,p_r$'' without
explicitly restricting $r$, so the singleton case, whose value is $p/(p-1)$,
must be excluded by hand; the modern restriction $|P|\ge2$ does that.

The letter's de-duplicated assertion is made for a general finite list of
primes and supplies no proof, so it is an asserted historical result and not an
open problem.  The note therefore restricts its open statement to the repeated
series $\mathcal R_P$.  A recovered proof of the letter's assertion, or a
transcendence statement for $\mathcal D_P$ with $|P|\ge3$, would be a different
question and needs its own formulation.

The two-prime argument of the note is independent of the unprinted argument in
the letter, and it is not the first public proof.  Steve Fan posted the same
factorisation, the same Hecke--Mahler reduction and the same conclusion in the
discussion thread of the problem's page on 26 June 2026
\cite{fan2026comment}, with follow-up remarks extending the argument to
arbitrary coprime pairs; a further comment there observes that the argument
does not seem to generalise immediately to $|P|\ge3$.  The manuscript of the
note was first released publicly on 22 July 2026, at commit \texttt{a9d3ab8}.

The statement of the problem has been formalised as a conjecture with an
unfilled proof in the \emph{Formal Conjectures} collection
\cite{formalconjectures269}.  Its Nat-indexed series includes the empty-prefix
least-common-multiple term, so its value is exactly \(1\) plus the conventional
series; transporting a theorem across that boundary needs an
explicit series-identification lemma, which is not supplied here.

\subsection{The abstract carry lemmas in their general form}
\label{long269:long:general-carries}

The note applies the carry machinery to the literal $\{2,3,5\}$ orbit.  The
underlying statements are about integer sequences alone, they are reusable, and
they are recorded here in the generality in which they are checked.

Let $D=D_{\mathrm{sm}}B$ with $D_{\mathrm{sm}}=2^{u}3^{v}5^{w}$ and
$\gcd(B,30)=1$, and let $(c_n)$ be an integer sequence satisfying
$c_{n+1}=b_nc_n-Dm_n$ for an integer radix word $(b_n)$ and forcing word
$(m_n)$.

\begin{proposition}[conditional denominator reduction]
\label{long269:long:denominator-reduction}
If $c_n=D_{\mathrm{sm}}d_n$ for every $n$, with $D_{\mathrm{sm}}>0$, then the
recurrence, positivity, upper bound and window identity for $(c_n)$ reduce to
the same four statements for $(d_n)$ with multiplier $B$ in place of $D$.
\end{proposition}

Height absorption alone does not imply divisibility of a carry state.  The
identity $DX_a=h_aN-Dv_a$ of Lemma~\ref{long269:res:all-scale-lattice} is what supplies
it for the actual orbit, and Theorem~\ref{long269:res:actual-cancellation} is the
instance in which the note uses it.  The formal consumer takes the
common-factor form as a hypothesis.

\begin{proposition}[conditional extinction of bounded carries]
\label{long269:long:windowconsumer}
Let $(b_n)$ and $(m_n)$ be a radix word and a forcing word, let
$G:\N_{>0}\times\N\to\N$, and assume cofinal local-window escape against $G$.
Fix $B>0$ coprime to $30$.  There is no integral sequence $(d_n)$ satisfying
simultaneously $d_{n+1}=b_nd_n-Bm_n$, $d_n>0$ and $|d_n|\le G(B,n)$ for every
$n\ge0$.
\end{proposition}

\begin{proof}
Choose one escaping window $(\ell,h)$.  The window identity gives
$d_{\ell+h}\equiv-BF_{\ell,h}$ modulo $|W_{\ell,h}|$.  The endpoint state is
positive and at most $G(B,\ell+h)$, whereas the canonical positive residue of
the right-hand side exceeds that bound, so Proposition~\ref{long269:res:consumer}
applies.
\end{proof}

Coprimality with $30$ is used by the escape hypothesis to select a window; once
a window is fixed, the finite contradiction does not use it.  The formalisation
carries the edge cases: a zero window base is excluded, a zero residue is
represented by the full modulus, and positivity prevents the endpoint carry
from vanishing.  The packaged absorbed form takes a nonzero smooth factor, the
exact factorisation $c_n=D_{\mathrm{sm}}d_n$, the absorbed recurrence for
$(c_n)$ and the positive short bound for $(d_n)$, and derives the same
contradiction in one statement.

\subsection{Worked finite examples}\label{long269:long:finite-geometry}

The ten smallest values of the running least common multiple at $\{2,3,5\}$ are
tabulated in Section~\ref{long269:sec:lcm}.  They illustrate both parts of
Proposition~\ref{long269:res:cell}: the value is constant on $\{5,6,7\}$ and on
$\{9,10\}$, and each change multiplies by a single prime, by $2$ at
$x=2,4,8$, by $3$ at $x=3,9$ and by $5$ at $x=5$.  Also $\Lprefix(10)=8\cdot9\cdot5=360$
is the least common multiple of the smooth numbers $1,2,3,4,5,6,8,9,10$.

For Proposition~\ref{long269:res:fibre} take $(p,q,r)=(2,3,5)$ and the box
$\mathcal B(1,1,1)$, whose eight points carry the smooth values
$1,2,3,5,6,10,15,30$ and the heights
\[
\begin{array}{c|cccccccc}
p^{i}q^{j}r^{k}&1&2&3&5&6&10&15&30\\ \hline
\hgt&1&2&6&60&60&360&360&10800
\end{array}
\]
Six heights occur, two of them twice: the points $5$ and $6$ share the height
$60$, and $10$ and $15$ share the height $360$.  The identity reads
\[
 1+\tfrac12+\tfrac16+\tfrac1{60}+\tfrac1{60}+\tfrac1{360}+\tfrac1{360}
 +\tfrac1{10800}
 =1+\tfrac12+\tfrac16+\tfrac2{60}+\tfrac2{360}+\tfrac1{10800}
 =\tfrac{18421}{10800},
\]
and the two coefficients $2$ carry the whole content of the regrouping on this
box.  Where the heights are pairwise distinct the identity is a relabelling.

Multiplying block radices along a run of blocks gives the product of the
jump-word letters over that run: for instance $b_1b_2=60$ is the product of the
four multipliers at $3,4,5,8$.  Block $4$ is the only one among the first six
carrying an internal power in both channels, and block $5$ carries neither, so
its radix falls back to the terminal factor alone.

\subsection{The finite-scan histogram}\label{long269:long:experiments}

The scan of Section~\ref{long269:sec:evidence} covers every $B\le5000$ coprime to $30$
and every start $100\le\ell\le3000$, with search depth permitted to $24$.  It
tested $3{,}869{,}934$ pairs, found an escaping window in every case, and the
distribution of first successful lengths was
\[
\resizebox{\linewidth}{!}{$\begin{array}{c|rrrrrrrrrrrrrrr}
h&4&5&6&7&8&9&10&11&12&13&14&15&16&17&18\\ \hline
\#&1&104&812&5437&51409&237423&735450&1431226&1132756&236752&34910&3076&521&49&8
\end{array}$}
\]
The first case at the maximal observed length $18$ is $B=917$ at start
$\ell=2980$, with endpoint jump index $6179$, window base
$18139852800000000$, forcing $13196471407660025821045$, residue
$76322101735$ and bound $3896420420$.  An earlier run of the same checker over
every $B\le1000$ coprime to $30$ and every start $100\le\ell\le500$ tested
$106{,}666$ pairs with maximal first successful length $14$, whose first case
is $B=359$ at start $291$, endpoint jump index $627$, base $5038848000000$,
forcing $25864575212865807$, residue $213175287$ and bound $15932659$.

By Corollary~\ref{long269:res:no-bounded-length} these histograms describe a bounded
region.  The observed concentration at lengths $10$ to $12$ sits just above the
necessary depth $\bigl(\log_2K(B,\ell+h)-\log_215\bigr)/3$ supplied by
Proposition~\ref{long269:res:window-growth}, which is what a reader should expect if
the residues behave like generic ones inside this range.

\subsection{Source inventory}\label{long269:long:sources}

Each entry names an informal statement and the Lean declaration that carries it
at the pinned source revision~\commitshort.

\begin{longtable}{@{}p{.35\linewidth}p{.59\linewidth}@{}}
informal statement & linked Lean source\\
\endhead
smooth lattice value & \lref{Erdos269/ThreePrimeRunningLcm.lean}{32}{smooth3Val}\\
pure-power height & \lref{Erdos269/ThreePrimeRunningLcm.lean}{37}{threePrimeHeight}\\
lattice kernel & \lref{Erdos269/ThreePrimeRunningLcm.lean}{41}{threePrimeKernelQ}\\
smooth prefix index set & \lref{Erdos269/ThreePrimeRunningLcm.lean}{47}{smoothPrefixExponents}\\
running least common multiple & \lref{Erdos269/ThreePrimeRunningLcm.lean}{54}{smoothPrefixLcm}\\
prefix value divides the height & \lref{Erdos269/ThreePrimeRunningLcm.lean}{60}{smooth3Val_dvd_threePrimeHeight_of_mem}\\
divisibility into the height & \lref{Erdos269/ThreePrimeRunningLcm.lean}{78}{smoothPrefixLcm_dvd_threePrimeHeight}\\
first pure-power membership & \lref{Erdos269/ThreePrimeRunningLcm.lean}{85}{pureFirst_mem_smoothPrefixExponents}\\
second pure-power membership & \lref{Erdos269/ThreePrimeRunningLcm.lean}{98}{pureSecond_mem_smoothPrefixExponents}\\
third pure-power membership & \lref{Erdos269/ThreePrimeRunningLcm.lean}{111}{pureThird_mem_smoothPrefixExponents}\\
running-lcm identity & \lref{Erdos269/ThreePrimeRunningLcm.lean}{124}{smoothPrefixLcm_eq_threePrimeHeight}\\
cell relation & \lref{Erdos269/ThreePrimeRunningLcm.lean}{164}{SameThreePrimeLogCell}\\
cell constancy of the height & \lref{Erdos269/ThreePrimeRunningLcm.lean}{170}{threePrimeHeight_eq_of_sameLogCell}\\
cell constancy of the running value & \lref{Erdos269/ThreePrimeRunningLcm.lean}{180}{smoothPrefixLcm_eq_of_sameLogCell}\\
cell constancy of the kernel & \lref{Erdos269/ThreePrimeRunningLcm.lean}{192}{threePrimeKernelQ_eq_of_sameLogCell}\\
positive power sets & \lref{Erdos269/ThreePrimeRunningLcm.lean}{205}{positivePrimePowers}\\
channel cardinality & \lref{Erdos269/ThreePrimeRunningLcm.lean}{209}{positivePrimePowers_card}\\
exclusion of the origin & \lref{Erdos269/ThreePrimeRunningLcm.lean}{220}{one_not_mem_positivePrimePowers}\\
channel disjointness & \lref{Erdos269/ThreePrimeRunningLcm.lean}{230}{positivePrimePowers_disjoint}\\
positive jump channels & \lref{Erdos269/ThreePrimeRunningLcm.lean}{244}{threePrimePositiveJumpSet}\\
positive jump count & \lref{Erdos269/ThreePrimeRunningLcm.lean}{250}{threePrimePositiveJumpSet_card}\\
jump set with the origin & \lref{Erdos269/ThreePrimeRunningLcm.lean}{274}{threePrimeJumpSetWithOrigin}\\
jump count with the origin & \lref{Erdos269/ThreePrimeRunningLcm.lean}{279}{threePrimeJumpSetWithOrigin_card}\\
first height step & \lref{Erdos269/ThreePrimeRunningLcm.lean}{295}{threePrimeHeight_firstLogStep}\\
second height step & \lref{Erdos269/ThreePrimeRunningLcm.lean}{306}{threePrimeHeight_secondLogStep}\\
third height step & \lref{Erdos269/ThreePrimeRunningLcm.lean}{317}{threePrimeHeight_thirdLogStep}\\
first coordinate step & \lref{Erdos269/ThreePrimeRunningLcm.lean}{327}{smoothPrefixLcm_firstLogStep}\\
second coordinate step & \lref{Erdos269/ThreePrimeRunningLcm.lean}{340}{smoothPrefixLcm_secondLogStep}\\
third coordinate step & \lref{Erdos269/ThreePrimeRunningLcm.lean}{353}{smoothPrefixLcm_thirdLogStep}\\
exponent box & \lref{Erdos269/ThreePrimeRunningLcm.lean}{369}{smoothExponentBox}\\
point height & \lref{Erdos269/ThreePrimeRunningLcm.lean}{374}{smoothPointHeight}\\
height fibre & \lref{Erdos269/ThreePrimeRunningLcm.lean}{378}{smoothHeightFiber}\\
fibre sum & \lref{Erdos269/ThreePrimeRunningLcm.lean}{385}{smoothHeightFiber_kernel_sum}\\
height-fibre normal form & \lref{Erdos269/ThreePrimeRunningLcm.lean}{407}{finiteSmoothKernelSum_groupedByHeight}\\
cubic majorant & \lref{Erdos269/ThreePrimeRunningLcm.lean}{431}{threePrimeHeight_le_cube}\\
origin value & \lref{Erdos269/ThreePrimeRunningLcm.lean}{444}{kernel_235_origin}\\
value at two & \lref{Erdos269/ThreePrimeRunningLcm.lean}{451}{kernel_235_two}\\
value at three & \lref{Erdos269/ThreePrimeRunningLcm.lean}{459}{kernel_235_three}\\
value at six & \lref{Erdos269/ThreePrimeRunningLcm.lean}{468}{kernel_235_six}\\
non-separation witness & \lref{Erdos269/ThreePrimeRunningLcm.lean}{480}{kernel_235_not_rankOne}\\
variable-base tail step & \lref{Erdos269/ThreePrimeRunningLcm.lean}{488}{tailStateStep}\\
expanded tail step & \lref{Erdos269/ThreePrimeRunningLcm.lean}{491}{tailStateStep_eq}\\
smooth exponent shell & \lref{Erdos269/ThreePrimeRunningLcm.lean}{501}{smoothExponentShell}\\
short-interval uniqueness & \lref{Erdos269/ThreePrimeRunningLcm.lean}{510}{exponent_unique_in_short_interval}\\
first-coordinate projection & \lref{Erdos269/ThreePrimeRunningLcm.lean}{544}{smoothExponentShell_card_le_dropFirst}\\
third-coordinate projection & \lref{Erdos269/ThreePrimeRunningLcm.lean}{586}{smoothExponentShell_card_le_dropThird}\\
sorted quadratic estimate & \lref{Erdos269/ThreePrimeRunningLcm.lean}{630}{sorted_pair_quadratic}\\
quadratic shell bound & \lref{Erdos269/ThreePrimeRunningLcm.lean}{647}{smoothExponentShell_card_quadratic}\\
dyadic internal power & \lref{Erdos269/ThreePrimeRunningLcm.lean}{663}{DyadicInternalPower}\\
internal-power uniqueness & \lref{Erdos269/ThreePrimeRunningLcm.lean}{669}{dyadicInternalPower_exponent_unique}\\
dyadic block base & \lref{Erdos269/ThreePrimeRunningLcm.lean}{751}{dyadicBlockBase235}\\
exact dyadic radix alphabet & \lref{Erdos269/ThreePrimeRunningLcm.lean}{762}{dyadicBlockBase235_cases}\\
bounded-radix consequence & \lref{Erdos269/ThreePrimeRunningLcm.lean}{774}{dyadicBlockBase235_mem_interval}\\
rank-two certificate & \lref{Erdos269/ThreePrimeRunningLcm.lean}{827}{kernel_235_minor_eq_neg_one_fifteen}\\
no integer rotation orbit & \lref{Erdos269/KernelCarryRank.lean}{154}{noIntegerOrbit_logb_of_prime}\\
height factorisation & \lref{Erdos269/KernelCarryRank.lean}{218}{threePrimeHeight_factorisation}\\
kernel factorisation & \lref{Erdos269/KernelCarryRank.lean}{259}{threePrimeKernelQ_factorisation}\\
two-prime outer product & \lref{Erdos269/KernelCarryRank.lean}{278}{twoPrimeKernelQ_eq_outer_product}\\
two-prime vanishing minors & \lref{Erdos269/KernelCarryRank.lean}{298}{twoPrimeKernelQ_minor_two_eq_zero}\\
uniform minors, general form & \lref{Erdos269/KernelCarryRank.lean}{313}{exists_uniform_nonsingular_threePrimeKernel_minor}\\
uniform minors for primes & \lref{Erdos269/KernelCarryRank.lean}{376}{exists_uniform_nonsingular_threePrimeKernel_minor_of_prime}\\
no finite separation & \lref{Erdos269/KernelCarryRank.lean}{388}{not_finite_separable_threePrimeKernel}\\
second-order minor & \lref{Erdos269/KernelCarryRank.lean}{560}{kernel_235_minor2_eq_neg_one_fifteen}\\
third-order minor & \lref{Erdos269/KernelCarryRank.lean}{570}{kernel_235_minor3_eq_one_over_81000}\\
misleading proportional row & \lref{Erdos269/KernelCarryRank.lean}{583}{kernel_235_row_three_eq_smul_row_zero}\\
failure of that proportionality & \lref{Erdos269/KernelCarryRank.lean}{592}{kernel_235_row_three_ne_smul_row_zero_at_four}\\
ordered block digit & \lref{Erdos269/DyadicBlockThresholdPartition.lean}{150}{dyadicOrderedBlockDigit235}\\
normalised state step & \lref{Erdos269/DyadicOrderedTailRecurrence.lean}{110}{dyadicNormalizedTailStateR235_succ}\\
shell summability & \lref{Erdos269/DyadicShellSummability.lean}{142}{summable_dyadicShellMassR235}\\
infinite tail recurrence & \lref{Erdos269/DyadicShellSummability.lean}{177}{dyadicNormalizedShellTsumTailR235_succ}\\
integer-or-cofinally-far dichotomy & \lref{Erdos269/DyadicShellSummability.lean}{188}{dyadicShellTsumTail_integer_or_cofinal_far}\\
boundary divisibility & \lref{Erdos269/RationalLatticeReduction.lean}{51}{smoothHeight_mul_prime_dvd_boundaryHeight}\\
half-height identity & \lref{Erdos269/RationalLatticeReduction.lean}{101}{two_mul_heightNormalizer235}\\
window clearing identity & \lref{Erdos269/RationalLatticeReduction.lean}{158}{heightNormalizer235_mul_windowMass_eq_int}\\
all-scale rationality lattice & \lref{Erdos269/RationalLatticeReduction.lean}{196}{qsmul_normalizedTailState_eq_int_of_value_eq_rat}\\
normalised-state collision & \lref{Erdos269/RationalLatticeReduction.lean}{296}{exists_normalizedTailState_collision_of_value_eq_rat}\\
least positive residue & \lref{Erdos269/ResidueEscape.lean}{26}{leastPositiveResidue}\\
positive representative range & \lref{Erdos269/ResidueEscape.lean}{31}{leastPositiveResidue_pos_le}\\
representative congruence & \lref{Erdos269/ResidueEscape.lean}{52}{leastPositiveResidue_modEq}\\
escape condition & \lref{Erdos269/ResidueEscape.lean}{71}{ResidueEscapesWindow}\\
finite natural-state contradiction & \lref{Erdos269/ResidueEscape.lean}{76}{no_bounded_positive_state_of_residue_escape}\\
contrapositive form & \lref{Erdos269/ResidueEscape.lean}{96}{residue_le_bound_of_bounded_positive_state}\\
integer least-positive-residue obstruction & \lref{Erdos269/ResidueEscape.lean}{110}{no_bounded_positive_int_state_of_leastPositiveResidue}\\
exact residue classifier & \lref{Erdos269/ResidueEscape.lean}{138}{leastPositiveResidue_eq_natAbs_of_pos_le_modEq}\\
window base & \lref{Erdos269/RestrictedFloorSum.lean}{417}{windowBase}\\
window forcing & \lref{Erdos269/RestrictedFloorSum.lean}{422}{windowForcing}\\
affine window identity & \lref{Erdos269/RestrictedFloorSum.lean}{455}{affineRecurrence_window}\\
scaled forcing identity & \lref{Erdos269/RestrictedFloorSum.lean}{469}{windowForcing_const_mul}\\
integral carry window & \lref{Erdos269/RestrictedFloorSum.lean}{480}{integralCarry_window}\\
endpoint residue identification & \lref{Erdos269/RestrictedFloorSum.lean}{497}{leastPositiveResidue_windowForcing_eq_carry}\\
absorbed smooth factor & \lref{Erdos269/RestrictedFloorSum.lean}{548}{smooth3Val_dvd_threePrimeHeight_of_le}\\
common-factor cancellation & \lref{Erdos269/RestrictedFloorSum.lean}{581}{integralCarry_cancel_commonFactor}\\
reduced bound transfer & \lref{Erdos269/RestrictedFloorSum.lean}{601}{reducedCarry_pos_le_of_commonFactor}\\
reduced window transfer & \lref{Erdos269/RestrictedFloorSum.lean}{612}{reducedIntegralCarry_window}\\
cofinal window hypothesis & \lref{Erdos269/RestrictedFloorSum.lean}{629}{CofinalLocalWindowEscape}\\
reduced-carry extinction & \lref{Erdos269/RestrictedFloorSum.lean}{645}{no_positive_reducedCarry_of_cofinalLocalWindowEscape}\\
absorbed-carry extinction & \lref{Erdos269/RestrictedFloorSum.lean}{689}{no_positive_absorbedCarry_of_cofinalLocalWindowEscape}\\
real window identity & \lref{Erdos269/CofinalWindowEscapeEquivalence.lean}{68}{trueNormalizedState_window}\\
escape from irrationality, general bound & \lref{Erdos269/CofinalWindowEscapeEquivalence.lean}{327}{cofinalLocalWindowEscape_of_irrational_of_beaten}\\
escape from irrationality, quadratic family & \lref{Erdos269/CofinalWindowEscapeEquivalence.lean}{355}{cofinalLocalWindowEscape_of_irrational_of_quadratic}\\
escape from irrationality, actual bound & \lref{Erdos269/CofinalWindowEscapeEquivalence.lean}{379}{cofinalLocalWindowEscape_of_irrational}\\
producer equivalence, shifted tail & \lref{Erdos269/CofinalWindowEscapeEquivalence.lean}{392}{actualCofinalLocalWindowEscape_iff}\\
producer equivalence, series value & \lref{Erdos269/CofinalWindowEscapeEquivalence.lean}{398}{actualCofinalLocalWindowEscape_iff_irrational_value}\\
rationality-to-carry bridge & \lref{Erdos269/RationalityCarryBridge.lean}{324}{exists_reducedCarry_of_value_eq_rat}\\
the actual producer & \lref{Erdos269/RationalityCarryBridge.lean}{397}{ActualCofinalLocalWindowEscape}\\
irrationality from the producer & \lref{Erdos269/RationalityCarryBridge.lean}{479}{irrational_value_of_cofinalLocalWindowEscape}\\
bounded-radix alternative & \lref{Erdos269/BoundedRadixTailEscape.lean}{89}{boundedRadix_zero_or_cofinal_far}\\
rational value from an integral scaled tail & \lref{Erdos269/BoundedRadixTailEscape.lean}{183}{rational_of_scaledTail_integer}\\
channel potential classifier & \lref{Erdos269/ThreeChannelBlockRigidity.lean}{59}{channelBlockNull_iff_channelPotential}\\
one-index lift extinction & \lref{Erdos269/CarryLiftExtinction.lean}{178}{no_carryLift_of_errorBound_below_twoPow}\\
uniform lift extinction & \lref{Erdos269/CarryLiftExtinction.lean}{238}{no_bounded_carryLift_of_blockNull_twoAnchors}\\
first-block sum & \lref{Erdos269/CarryLiftExtinction.lean}{289}{binaryLift_twoAnchors_force_firstBlock_sum_one}\\
four-state obstruction & \lref{Erdos269/CarryLiftExtinction.lean}{308}{no_unitAccurateLift_with_twoAnchors_and_firstTwoBlockNull}\\
residue-coboundary form & \lref{Erdos269/WeightedPhaseCarry.lean}{109}{carry_eq_residueDigit_add_coboundary}\\
carry interval & \lref{Erdos269/WeightedPhaseCarry.lean}{150}{carryResidue_mem_interval}\\
digit interval & \lref{Erdos269/WeightedPhaseCarry.lean}{157}{residueDigit_mem_interval}\\
symbolic realisation & \lref{Erdos269/WeightedPhaseCarry.lean}{293}{CartierRealisation}\\
finite realised span & \lref{Erdos269/WeightedPhaseCarry.lean}{334}{finite_realisedSpan_of_factorisation}\\
Cantor-state confinement & \lref{Erdos269/PurePowerIrrationality.lean}{61}{state_mem_Ioo}\\
state gap bound & \lref{Erdos269/PurePowerIrrationality.lean}{74}{stateGap_lower_bound}\\
small-form criterion & \lref{Erdos269/PurePowerIrrationality.lean}{94}{irrational_of_small_nonzero_forms}\\
conditional assembly & \lref{Erdos269/PurePowerIrrationality.lean}{156}{irrational_of_clearing_and_small_gaps}\\
staircase index engine & \lword{Shared/IrrationalRotationStaircase.lean}{271}{exists_staircase_indices}{staircase indices}\\
\end{longtable}
